% !TEX program = xelatex
%%%%%%%%%%%%%%%%%%%%%%%%%%%%%%%%%%%%%%%%%%%%%%%%%%%%%%%%%%%%%%%%%%%%%%%%%%%%%%%
%  SCX-Audited Blockchain Consensus:
%  Multi-Validator Certification of Distributed Ledger Integrity
%  SCX Research Collective — June 2026
%%%%%%%%%%%%%%%%%%%%%%%%%%%%%%%%%%%%%%%%%%%%%%%%%%%%%%%%%%%%%%%%%%%%%%%%%%%%%%%
\documentclass[11pt,a4paper]{article}

% ── SCX Framework Standard Packages ──────────────────────────
\usepackage{fontspec}
\usepackage{amsmath,amssymb,amsthm}
\usepackage{mathtools}
\usepackage{bm}
\usepackage{graphicx}
\usepackage{xcolor}
\usepackage{hyperref}
\usepackage{geometry}
\usepackage{enumitem}
\usepackage{booktabs}
\usepackage{multirow}
\usepackage{array}
\usepackage{algorithm}
\usepackage{algpseudocode}

% CJK support via fontspec (XeLaTeX/LuaLaTeX)
\usepackage{xeCJK}
\setCJKmainfont{SimSun}

\geometry{margin=1in}

% ── SCX Theorems & Proofs ─────────────────────────────────────
\newtheorem{theorem}{Theorem}[section]
\newtheorem{proposition}[theorem]{Proposition}
\newtheorem{lemma}[theorem]{Lemma}
\newtheorem{corollary}[theorem]{Corollary}
\newtheorem{conjecture}[theorem]{Conjecture}
\newtheorem{definition}[theorem]{Definition}
\newtheorem{remark}[theorem]{Remark}
\theoremstyle{definition}
\newtheorem{assumption}[theorem]{Assumption}
\newtheorem{example}[theorem]{Example}

% ── Rigor Labels ──────────────────────────────────────────────
\newcommand{\rigorFull}[0]{$\blacksquare$~\textbf{[Full Proof]}}
\newcommand{\rigorPartial}[0]{$\square$~\textbf{[Partial Proof]}}
\newcommand{\rigorSketch}[0]{$\diamond$~\textbf{[Proof Sketch]}}

% ── SCX Framework Macros ──────────────────────────────────────
\newcommand{\Situs}{\textsc{Situs}}
\newcommand{\Yajie}{\textsc{Yajie}}
\newcommand{\Spring}{\textsc{Spring}}
\newcommand{\Cercis}{\textsc{Cercis}}
\newcommand{\SCX}{\textsc{SCX}}

% ── Math Operators ────────────────────────────────────────────
\DeclareMathOperator*{\argmin}{arg\,min}
\DeclareMathOperator*{\argmax}{arg\,max}
\DeclareMathOperator{\TV}{TV}
\DeclareMathOperator{\KL}{KL}
\DeclareMathOperator{\Var}{Var}
\DeclareMathOperator{\Cov}{Cov}
\DeclareMathOperator{\Bias}{Bias}
\DeclareMathOperator{\supp}{supp}
\DeclareMathOperator{\Lip}{Lip}
\DeclareMathOperator{\SHA}{SHA}

% ── Notation ──────────────────────────────────────────────────
\newcommand{\R}{\mathbb{R}}
\newcommand{\N}{\mathbb{N}}
\newcommand{\E}{\mathbb{E}}
\newcommand{\Pbb}{\mathbb{P}}
\newcommand{\ind}[1]{\mathbf{1}[#1]}
\newcommand{\norm}[1]{\left\lVert#1\right\rVert}
\newcommand{\abs}[1]{\left|#1\right|}
\newcommand{\inner}[2]{\langle #1, #2 \rangle}
\newcommand{\red}{\textcolor{red}}
\newcommand{\cB}{\mathcal{B}}
\newcommand{\cC}{\mathcal{C}}
\newcommand{\cD}{\mathcal{D}}
\newcommand{\cH}{\mathcal{H}}
\newcommand{\cS}{\mathcal{S}}
\newcommand{\cV}{\mathcal{V}}
\newcommand{\cM}{\mathcal{M}}

% ── Blockchain Domain Macros ──────────────────────────────────
\newcommand{\chain}{\mathsf{Chain}}
\newcommand{\block}{\mathsf{Block}}
\newcommand{\state}{\mathsf{State}}
\newcommand{\tx}{\mathsf{tx}}
\newcommand{\valSet}{\mathcal{V}}
\newcommand{\honestSet}{\mathcal{H}}
\newcommand{\advSet}{\mathcal{A}}
\newcommand{\stake}{\mathsf{stake}}
\newcommand{\slashPenalty}{\kappa}
\newcommand{\Meff}{M_{\mathsf{eff}}}
\newcommand{\Mthresh}{M_{\mathsf{thresh}}}
\newcommand{\Mt}{M_t}
\newcommand{\consensusState}{c}

% ── Assumption / Limitation Tags ──────────────────────────────
\newcommand{\assumptionTag}[1]{\textsf{\textbf{[A#1]}}}
\newcommand{\limitationTag}[1]{\textsf{\textbf{[L#1]}}}

% ── Info ──────────────────────────────────────────────────────
\title{\textbf{SCX-Audited Blockchain Consensus:}\\
\large Multi-Validator Certification of Distributed Ledger Integrity\\[4pt]
\small SCX审计区块链共识：分布式账本完整性的多验证者认证}

\author{SCX Research Collective}
\date{\today}

\begin{document}
\maketitle

\begin{abstract}
We formalize blockchain consensus as a multi-expert certification problem under the \SCX{} (Structured Causal eXamination) auditing framework. Each block is modeled as a \textbf{claim about the ledger state}, and validators are treated as \textbf{multi-expert auditors} producing independent certifications. We prove three core theorems. \textbf{Theorem~1 (Fork Boundary Theorem 分叉界定理):} The probability of a consensus fork—where two conflicting blocks are certified at the same height—is bounded by $\exp(-2M\varepsilon^2)$ for $M$ independent validators with minimum detection gap $\varepsilon$, establishing that forks represent consensus failure with exponentially decaying likelihood as validator multiplicity grows. \textbf{Theorem~2 (51\% Attack as Effective-M Collapse 51\%攻击作为有效M塌缩):} When an adversary controls fraction $\alpha > 0.5$ of validation power, the effective auditor multiplicity $\Meff$ drops below the consensus threshold $\Mthresh = 1$, formally proving that 51\% attacks are equivalent to auditor capture collapsing the certification guarantee. \textbf{Theorem~3 (PoS Slashing as Audit Penalty 权益证明罚没作为审计惩罚定理):} In proof-of-stake systems, the slashing penalty $\slashPenalty$ is structurally identical to the \SCX{} audit penalty $\kappa$, and we establish the \Yajie{} Nash-Pareto Equilibrium condition $\slashPenalty > \slashPenalty^*$ under which honest validation strictly dominates equivocation. We further demonstrate that the hash chain implements \Spring's permanent memory $\Mt$—the auditor multiplicity at each block height is permanently and immutably recorded—and that Nakamoto consensus constitutes an \textbf{implicit M-declaration without formal guarantee}, highlighting the structural gap that \SCX{} fills with explicit, provable certification. The framework is illustrated through formal analysis of Bitcoin, Ethereum PoS, Tendermint, and Solana consensus mechanisms.

\medskip
\noindent\textbf{Keywords:} SCX auditing, blockchain consensus 区块链共识, multi-validator certification 多验证者认证, fork boundary 分叉界, 51\% attack 51\%攻击, proof-of-stake slashing 权益证明罚没, \Spring{} permanent memory, Nakamoto consensus 中本聪共识, distributed ledger integrity 分布式账本完整性
\end{abstract}

%=====================================================================
\section{Introduction 引言}
%=====================================================================

Blockchain systems achieve distributed consensus without a central authority: a network of validators (miners in proof-of-work, stakers in proof-of-stake) collectively certifies the sequence of state transitions that constitute the ledger. The security guarantee is probabilistic: as long as a supermajority of validators follows the protocol, the ledger remains consistent, immutable, and resistant to adversarial manipulation. This guarantee, however, is typically stated in aggregate terms—``the network is secure if 51\% of hash power is honest''—without formalizing \textit{what it means for each individual block to be certified} and \textit{how many independent validators actually examined it}.

The \SCX{} auditing framework~\cite{SCX2025} provides precisely the missing formalism. \SCX{} models verification as a multi-expert signaling game: an entity publishes a \textbf{claim} about a state, and $M$ independent auditors each produce estimates of the true state. The consensus among auditors, quantified by the \Yajie{} Nash-Pareto Equilibrium, provides an auditable confidence bound on the claim's veracity. The \Spring{} mechanism permanently records the audit history, and the \Cercis{} score combines quality and novelty into a single metric.

\textbf{The Central Analogy.} We observe a structural isomorphism between blockchain consensus and \SCX{} multi-expert audit:

\begin{enumerate}[label=(\roman*)]
    \item \textbf{Block $\leftrightarrow$ Claim.} Each block proposes a new ledger state—a claim that ``after executing transactions $T_1, \ldots, T_k$, the state is $S'$.'' This is structurally identical to a scientific claim ``after applying methodology $\mathcal{M}$ to data $\mathcal{D}$, the result is $R$.''
    \item \textbf{Validator $\leftrightarrow$ Auditor.} Each validator independently verifies the block's correctness: checking transaction signatures, state transition validity, and consensus rules. This mirrors an auditor independently estimating the true state from raw data.
    \item \textbf{Consensus $\leftrightarrow$ \Yajie{} Agreement.} When $M$ validators certify the same block, their agreement constitutes a \Yajie{} consensus with detection power proportional to $M$.
    \item \textbf{Fork $\leftrightarrow$ Consensus Failure.} A blockchain fork—where two conflicting blocks gain certification at the same height—is structurally identical to \SCX{} audit failure: the multi-expert panel disagrees on the true state.
    \item \textbf{51\% Attack $\leftrightarrow$ Auditor Capture.} When an adversary controls a majority of validators, the effective auditor multiplicity $\Meff$ collapses below the threshold required for reliable certification.
    \item \textbf{Slashing $\leftrightarrow$ Audit Penalty $\kappa$.} In proof-of-stake, validators who certify conflicting blocks lose their stake. This is exactly the \SCX{} audit penalty: the cost imposed when a claim is detected as false by the multi-expert panel.
\end{enumerate}

\textbf{Contributions.} This paper provides:

\begin{enumerate}[label=(\roman*)]
    \item \textbf{Formalization} (Section~\ref{sec:formal}): Blockchain consensus as an \SCX{} multi-expert certification game with explicit assumptions.
    \item \textbf{Three theorems with full proofs}:
    \begin{itemize}
        \item Theorem~\ref{thm:fork}: Fork Boundary Theorem—exponential bound on fork probability under $M$ validators.
        \item Theorem~\ref{thm:fiftyone}: 51\% Attack as Effective-M Collapse—formal equivalence between adversarial majority and auditor capture.
        \item Theorem~\ref{thm:slashing}: PoS Slashing as Audit Penalty—\Yajie{} equilibrium condition for honest validation dominance.
    \end{itemize}
    \item \textbf{\Spring{} Permanent Memory} (Section~\ref{sec:spring}): The hash chain as \Spring's permanent record $\Mt$, providing immutable audit trails at every block height.
    \item \textbf{Nakamoto Consensus Critique} (Section~\ref{sec:nakamoto}): Nakamoto consensus as implicit M-declaration without formal guarantee—why explicit $M$ matters.
    \item \textbf{Consensus Mechanism Analysis} (Section~\ref{sec:mechanisms}): Application of the framework to Bitcoin, Ethereum PoS, Tendermint, and Solana.
    \item \textbf{Discussion} (Section~\ref{sec:discussion}): Limitations, extensions, and the path toward formally audited blockchain consensus.
\end{enumerate}

\textbf{What this paper is not.} This is not a new consensus protocol proposal. We do not claim to replace Nakamoto consensus, Tendermint, or any existing mechanism. Rather, we provide a \textit{verification framework}—a mathematical lens through which existing consensus mechanisms can be audited, their guarantees can be rigorously stated, and their failure modes can be precisely characterized. The value is in making implicit assumptions explicit and in providing formal tools for reasoning about consensus security.

%=====================================================================
\section{The SCX Blockchain Audit Framework 区块链审计框架}
\label{sec:framework}
%=====================================================================

\subsection{Blocks as Claims About State 区块作为状态声明}

\begin{definition}[Ledger State 账本状态]
\label{def:state}
At any block height $h$, the ledger state $S_h$ is a finite mapping from accounts to balances (for cryptocurrency) or, more generally, a finite mapping from addresses to arbitrary state variables (for smart contract platforms):
\begin{equation}
    S_h: \mathcal{A} \to \mathcal{V},
    \label{eq:state}
\end{equation}
where $\mathcal{A}$ is the address space and $\mathcal{V}$ is the value space. For Bitcoin, $\mathcal{V} \subset \R_{\geq 0}$ (UTXO set). For Ethereum, $\mathcal{V}$ includes account nonces, balances, storage roots, and code hashes.
\end{definition}

\begin{definition}[Block as State Claim 区块作为状态声明]
\label{def:block_claim}
A block $B_h = (H, T, \sigma)$ at height $h$ is a \textbf{claim} that:
\begin{equation}
    \text{``Starting from state } S_{h-1}, \text{ after executing transactions } T = \{tx_1, \ldots, tx_k\}, \text{ the new state is } S_h.''
    \label{eq:block_claim}
\end{equation}
The block contains: (i)~a header $H$ with parent hash, timestamp, and Merkle root; (ii)~the transaction list $T$; (iii)~cryptographic signatures $\sigma$ from validators who certify the block. The claim is \textbf{verifiable}: any party can re-execute $T$ from $S_{h-1}$ and check whether the resulting state matches $S_h$.
\end{definition}

\begin{definition}[Blockchain as Claim Chain 区块链作为声明链]
\label{def:blockchain}
A blockchain $\chain = (B_0, B_1, \ldots, B_H)$ is a sequence of state claims where each claim $B_h$ references its predecessor $B_{h-1}$ via cryptographic hash:
\begin{equation}
    H.\text{parent\_hash} = \SHA256(B_{h-1}).
    \label{eq:hash_chain}
\end{equation}
The chain is \textbf{consistent} if for all $h \geq 1$, the state transition $S_{h-1} \xrightarrow{T_h} S_h$ is valid according to the protocol rules. The chain is \textbf{fully certified} if every block $B_h$ carries signatures from a set of validators $\valSet_h$ satisfying the protocol's certification threshold.
\end{definition}

\subsection{Validators as Multi-Expert Auditors 验证者作为多专家审计者}

\begin{definition}[Validator as Auditor 验证者作为审计者]
\label{def:validator}
A validator $V_j \in \valSet$ is an \textbf{auditor} who, upon receiving a block proposal $B_h$, performs:
\begin{enumerate}[label=(\roman*)]
    \item \textbf{State verification}: Re-execute all transactions $T$ against $S_{h-1}$ and confirm the resulting state matches $S_h$.
    \item \textbf{Signature verification}: Validate all cryptographic signatures in $T$.
    \item \textbf{Consensus rule check}: Verify that $B_h$ satisfies all protocol-level rules (block size, gas limits, timestamp validity, etc.).
    \item \textbf{Certification}: If all checks pass, produce a signature $\sigma_j$ attesting to $B_h$'s validity.
\end{enumerate}
The validator outputs an estimate of the block's validity: $\hat{v}_j(B_h) \in \{0, 1\}$, where $1$ indicates certification and $0$ indicates rejection. The estimates are independent conditional on the true validity $v(B_h) \in \{0, 1\}$.
\end{definition}

\begin{definition}[Validator Multiplicity $M$ 验证者多重性]
\label{def:M}
For a block $B_h$, the \textbf{validator multiplicity} $M_h = |\valSet_h|$ is the number of validators who independently certify the block. The \textbf{effective multiplicity} $\Meff$ accounts for validator correlation:
\begin{equation}
    \Meff = \frac{M}{1 + (M-1)\bar{\rho}},
    \label{eq:Meff}
\end{equation}
where $\bar{\rho} \in [0, 1]$ is the average pairwise correlation of validator errors. Independent validators ($\bar{\rho} = 0$) yield $\Meff = M$; perfectly correlated validators ($\bar{\rho} = 1$) yield $\Meff = 1$.
\end{definition}

\begin{remark}[M-Declaration in Blockchain]
In the \SCX{} framework, every claim must declare its $M$ parameter. For blockchain, the natural $M$-declaration is the number of validators who certified each block—which is already recorded in block headers (e.g., validator signatures in Tendermint, attestations in Ethereum PoS). The blockchain provides \textbf{automatic M-declaration}: you cannot certify a block without recording who certified it. This is a structural advantage of blockchain over traditional scientific publication, where $M$ must be explicitly declared and can be omitted.
\end{remark}

\subsection{Consensus and Fork as Audit Outcome 共识与分叉作为审计结果}

\begin{definition}[Consensus 共识]
\label{def:consensus}
A block $B_h$ achieves \textbf{consensus} at multiplicity $M$ if at least $M$ validators independently certify $B_h$ and no other block at height $h$ receives $\geq M$ certifications. The \Yajie{} consensus value is:
\begin{equation}
    \consensusState(B_h) = \ind{\sum_{j=1}^{M} \ind{\hat{v}_j(B_h) = 1} \geq \tau M},
    \label{eq:consensus}
\end{equation}
where $\tau \in (0.5, 1]$ is the consensus threshold (typically $\tau = 2/3$ for BFT protocols).
\end{definition}

\begin{definition}[Fork as Consensus Failure 分叉作为共识失败]
\label{def:fork}
A \textbf{fork} at height $h$ occurs when two distinct blocks $B_h \neq B_h'$ both achieve consensus:
\begin{equation}
    \consensusState(B_h) = \consensusState(B_h') = 1.
    \label{eq:fork}
\end{equation}
In \SCX{} terms, a fork is an \textbf{audit failure}: the multi-expert panel of validators disagrees on the true state. This is structurally identical to a scientific audit where the expert panel splits on whether a claim is valid—the \SCX{} framework provides the mathematical tools to bound the probability and severity of such failures.
\end{definition}

%=====================================================================
\section{Formal Model 形式模型}
\label{sec:formal}
%=====================================================================

\subsection{The Consensus Game 共识博弈}

We model blockchain consensus as a multi-player certification game. The players are $N$ validators $\valSet = \{V_1, \ldots, V_N\}$. Nature selects the true state transition validity $v \in \{0, 1\}$, where $v = 1$ means the proposed block is valid and $v = 0$ means it is invalid (e.g., contains double-spends).

\begin{definition}[Validator Payoff 验证者收益]
\label{def:payoff}
Validator $V_j$ receives:
\begin{itemize}
    \item \textbf{Reward} $R > 0$ for certifying a block that achieves consensus (regardless of validity).
    \item \textbf{Penalty} $\slashPenalty > 0$ for certifying an invalid block that is later detected as invalid (slashing in PoS, wasted work in PoW).
    \item \textbf{No reward or penalty} for rejecting a block or for certifying a block that does not achieve consensus.
\end{itemize}
The expected payoff for action $a_j \in \{0, 1\}$ (reject/certify) is:
\begin{equation}
    U_j(a_j; v) = a_j \cdot \big[R \cdot \Pbb(\text{consensus} \mid a_j) - \slashPenalty \cdot (1 - v) \cdot \Pbb(\text{detection} \mid v=0)\big].
    \label{eq:validator_payoff}
\end{equation}
\end{definition}

\begin{definition}[Adversarial Validator 对抗性验证者]
\label{def:adversary}
An adversary $\advSet \subset \valSet$ controls fraction $\alpha = |\advSet| / N$ of validators. Adversarial validators may deviate arbitrarily from the protocol: certifying invalid blocks, withholding certifications, or equivocating (certifying conflicting blocks). Honest validators $\honestSet = \valSet \setminus \advSet$ follow the protocol exactly: certify valid blocks, reject invalid ones, never equivocate.
\end{definition}

\subsection{Assumptions 假设}

We now state the assumptions under which our theorems hold.

\begin{assumption}[A1: Bounded Validator Set 有界验证者集合]
\label{ass:A1}
The validator set $\valSet$ has finite size $N = |\valSet| < \infty$. This holds for any practical blockchain: Bitcoin has a finite number of mining pools, Ethereum PoS has a finite validator registry, and Tendermint has a fixed validator set per epoch.
\end{assumption}

\begin{assumption}[A2: Independent Validation 独立验证]
\label{ass:A2}
Conditional on the true block validity $v$, honest validators' certification decisions are mutually independent:
\begin{equation}
    \Pbb(\hat{v}_1, \ldots, \hat{v}_{|\honestSet|} \mid v) = \prod_{j \in \honestSet} \Pbb(\hat{v}_j \mid v).
    \label{eq:independence}
\end{equation}
This holds when validators use independent hardware, network paths, and software implementations. Correlation may arise from shared software bugs or network partitions; these are captured by the effective multiplicity adjustment in Eq.~\eqref{eq:Meff}.
\end{assumption}

\begin{assumption}[A3: Bounded Error Rate 有界错误率]
\label{ass:A3}
Each honest validator $V_j$ has error probability $p_j \in [0, 0.5)$:
\begin{equation}
    p_j = \Pbb(\hat{v}_j = 0 \mid v = 1) + \Pbb(\hat{v}_j = 1 \mid v = 0) \leq p_{\max} < 0.5.
    \label{eq:error_rate}
\end{equation}
The minimum detection gap is $\varepsilon = 0.5 - p_{\max} > 0$. Validators with $p_j \geq 0.5$ are indistinguishable from random guessing and provide no certification value.
\end{assumption}

\begin{assumption}[A4: Honest Supermajority (Conditional) 诚实超多数（有条件）]
\label{ass:A4}
For Theorems~1 and~3, we assume honest validators constitute a supermajority: $\alpha < 0.5$. Theorem~2 examines what happens when this assumption is violated ($\alpha \geq 0.5$).
\end{assumption}

\begin{assumption}[A5: Slashing Enforceability 罚没可执行性]
\label{ass:A5}
In proof-of-stake systems, equivocation (certifying conflicting blocks) is cryptographically detectable and the slashing penalty $\slashPenalty$ is automatically enforced by the protocol. This holds for Ethereum PoS, Tendermint, and all BFT-based PoS systems.
\end{assumption}

\begin{assumption}[A6: Hash Preimage Resistance 哈希原像抗性]
\label{ass:A6}
The hash function $\SHA256$ used for chain linking satisfies cryptographic preimage resistance: given $y = \SHA256(x)$, finding $x' \neq x$ with $\SHA256(x') = y$ is computationally infeasible.
\end{assumption}

\begin{assumption}[A7: Synchronous Network (Bounded Delay) 同步网络（有界延迟）]
\label{ass:A7}
Messages between honest validators are delivered within a known maximum delay $\Delta$. Validators have loosely synchronized clocks. This is the standard partial synchrony assumption for BFT consensus protocols.
\end{assumption}

%=====================================================================
\section{Theorem 1: Fork Boundary Theorem 分叉界定理}
\label{sec:theorem1}
%=====================================================================

We now prove that the probability of a consensus fork decays exponentially in the effective validator multiplicity $\Meff$, establishing a rigorous bound on consensus failure.

\subsection{Intuition 直观理解}

A fork occurs when the validator set splits its certifications between two conflicting blocks. For a fork to be ``meaningful''—i.e., to create a persistent chain split—both blocks must receive certifications from a substantial fraction of validators. If validators are independent and honest, the probability of such a split is bounded by the probability that the mean of independent Bernoulli trials deviates from its expectation by a large margin. Hoeffding's inequality provides the exponential bound.

In the \SCX{} framework, this is identical to the problem of \textit{auditor disagreement}: when $M$ experts independently evaluate a claim, what is the probability that they split into two camps of roughly equal size? The answer is that this probability vanishes exponentially as $M$ grows, provided the experts have non-trivial individual accuracy ($p_j < 0.5$).

\begin{lemma}[Certification Concentration 认证集中不等式]
\label{lem:cert_concentration}
Let $M$ honest validators independently certify a valid block ($v = 1$), each with error probability $p_j \leq p_{\max} < 0.5$. Let $X = \sum_{j=1}^{M} \ind{\hat{v}_j = 1}$ be the number of certifications. For any $\delta > 0$:
\begin{equation}
    \Pbb\left(\frac{X}{M} \leq 1 - p_{\max} - \delta\right) \leq \exp(-2M\delta^2).
    \label{eq:cert_concentration}
\end{equation}
Similarly, for an invalid block ($v = 0$):
\begin{equation}
    \Pbb\left(\frac{X}{M} \geq p_{\max} + \delta\right) \leq \exp(-2M\delta^2).
    \label{eq:cert_concentration_invalid}
\end{equation}
\end{lemma}

\begin{proof}\rigorFull
\textbf{Step 1: Valid block certification.} For $v = 1$, each validator certifies with probability $1 - p_j \geq 1 - p_{\max}$. Define $Y_j = 1 - \ind{\hat{v}_j = 1}$, so $\E[Y_j] = p_j \leq p_{\max}$. Then $X/M \leq 1 - p_{\max} - \delta$ is equivalent to $\frac{1}{M}\sum_{j=1}^{M} Y_j \geq p_{\max} + \delta$. By Hoeffding's inequality for bounded independent random variables $Y_j \in [0, 1]$:
\begin{equation}
    \Pbb\left(\frac{1}{M}\sum_{j=1}^{M} Y_j \geq \frac{1}{M}\sum_{j=1}^{M} \E[Y_j] + \delta\right) \leq \exp(-2M\delta^2).
\end{equation}
Since $\frac{1}{M}\sum \E[Y_j] \leq p_{\max}$, the event $\frac{1}{M}\sum Y_j \geq p_{\max} + \delta$ is contained in the Hoeffding event, yielding Eq.~\eqref{eq:cert_concentration}.

\textbf{Step 2: Invalid block certification.} For $v = 0$, each validator erroneously certifies with probability $p_j \leq p_{\max}$. Direct application of Hoeffding to $Z_j = \ind{\hat{v}_j = 1}$ with $\E[Z_j] = p_j \leq p_{\max}$ yields Eq.~\eqref{eq:cert_concentration_invalid}.

\textbf{Step 3: Correlation adjustment.} When validators are correlated with average pairwise correlation $\bar{\rho}$, we replace $M$ with $\Meff = M / (1 + (M-1)\bar{\rho})$ by the effective sample size formula for correlated binomial trials~\cite{hoeffding1963}. The bound becomes $\exp(-2\Meff\delta^2)$. $\square$
\end{proof}

\begin{theorem}[Fork Boundary Theorem 分叉界定理]
\label{thm:fork}
\rigorFull
Let there be $N$ validators, of which $M = |\honestSet| = (1-\alpha)N$ are honest with maximum error probability $p_{\max} < 0.5$. Under Assumptions~\ref{ass:A1}--\ref{ass:A4} and~\ref{ass:A7}, for a consensus threshold $\tau \in (p_{\max}, 1-p_{\max})$, the probability of a fork at height $h$ is bounded by:
\begin{equation}
    \Pbb(\text{fork}_h) \leq 2\exp\left(-2\Meff \cdot (\tau - p_{\max})^2\right),
    \label{eq:fork_bound}
\end{equation}
where $\Meff = M / (1 + (M-1)\bar{\rho})$ is the effective honest validator multiplicity. As $M \to \infty$, $\Pbb(\text{fork}_h) \to 0$ exponentially fast.

Furthermore, the expected number of forks over $H$ blocks satisfies:
\begin{equation}
    \E[\text{\#forks in } H \text{ blocks}] \leq 2H \cdot \exp\left(-2\Meff \cdot (\tau - p_{\max})^2\right).
    \label{eq:expected_forks}
\end{equation}
\end{theorem}

\begin{proof}\rigorFull
\textbf{Step 1: Fork condition.} A fork occurs when two distinct blocks $B_h \neq B_h'$ both receive $\geq \tau M$ certifications from honest validators. For a valid block $B_h$, the fraction of honest validators certifying is $X/M$. For an alternative block $B_h'$ (which must be invalid, since only one valid block can exist at height $h$ under deterministic state transitions), the fraction certifying is $X'/M$. A fork requires $X/M \geq \tau$ AND $X'/M \geq \tau$.

\textbf{Step 2: Valid block shortfall.} For the valid block $B_h$, the event $X/M < \tau$ (i.e., fewer than $\tau M$ honest validators certify) occurs when $X/M \leq 1 - p_{\max} - (\tau - p_{\max})$. By Lemma~\ref{lem:cert_concentration} with $\delta = \tau - p_{\max}$ (which is positive since $\tau > p_{\max}$ by assumption):
\begin{equation}
    \Pbb\left(\frac{X}{M} < \tau\right) \leq \exp\left(-2\Meff \cdot (\tau - p_{\max})^2\right).
    \label{eq:valid_shortfall}
\end{equation}

\textbf{Step 3: Invalid block certification.} For the invalid block $B_h'$, the event $X'/M \geq \tau$ requires at least $\tau M$ honest validators to erroneously certify an invalid block. For this to happen, $X'/M \geq p_{\max} + (\tau - p_{\max})$. By Lemma~\ref{lem:cert_concentration} with $\delta = \tau - p_{\max}$:
\begin{equation}
    \Pbb\left(\frac{X'}{M} \geq \tau\right) \leq \exp\left(-2\Meff \cdot (\tau - p_{\max})^2\right).
    \label{eq:invalid_cert}
\end{equation}

\textbf{Step 4: Union bound.} The fork event is contained in the union of the two events above (either the valid block fails to reach threshold, or an invalid block spuriously reaches threshold, or both). By union bound:
\begin{equation}
    \Pbb(\text{fork}_h) \leq \Pbb\left(\frac{X}{M} < \tau\right) + \Pbb\left(\frac{X'}{M} \geq \tau\right) \leq 2\exp\left(-2\Meff \cdot (\tau - p_{\max})^2\right).
\end{equation}

\textbf{Step 5: Expected forks over $H$ blocks.} By linearity of expectation and independence across heights (conditional on validator behavior being i.i.d.):
\begin{equation}
    \E[\text{\#forks in } H \text{ blocks}] = \sum_{h=1}^{H} \Pbb(\text{fork}_h) \leq 2H \cdot \exp\left(-2\Meff \cdot (\tau - p_{\max})^2\right).
\end{equation}

\textbf{Step 6: Asymptotic behavior.} As $M \to \infty$, $\Meff \to \infty$ (since $\Meff \geq M/(1+M\bar{\rho}) \to 1/\bar{\rho}$ for $\bar{\rho} > 0$, and $\Meff = M \to \infty$ for $\bar{\rho} = 0$), so $\exp(-2\Meff \cdot (\tau - p_{\max})^2) \to 0$. The fork probability vanishes exponentially. For any $\varepsilon > 0$, choosing $M > \frac{\ln(2/\varepsilon)}{2(\tau - p_{\max})^2}$ ensures $\Pbb(\text{fork}_h) < \varepsilon$. $\square$
\end{proof}

\begin{corollary}[Fork-Free Threshold 无分叉阈值]
\label{cor:fork_free}
For a target fork probability $\varepsilon$ (e.g., $\varepsilon = 10^{-9}$), the minimum honest validator multiplicity is:
\begin{equation}
    M_{\min}(\varepsilon) = \left\lceil \frac{(1 + (M-1)\bar{\rho}) \cdot \ln(2/\varepsilon)}{2(\tau - p_{\max})^2} \right\rceil.
    \label{eq:M_min_fork}
\end{equation}
For typical BFT parameters ($\tau = 2/3$, $p_{\max} = 10^{-3}$, $\bar{\rho} = 0.01$), this yields $M_{\min}(10^{-9}) \approx 47$ honest validators. For Nakamoto consensus with $\tau \approx 0.5$ and higher $p_{\max}$, $M_{\min}$ is substantially larger.
\end{corollary}

\begin{remark}[Forks = Consensus Failure in SCX]
Theorem~\ref{thm:fork} formalizes the \SCX{} principle that \textbf{forks are not a feature—they are a failure of multi-expert certification}. In the \SCX{} framework, when $M$ independent experts disagree on a claim's validity, the audit has failed to produce a definitive answer. The probability of this failure is bounded and controllable: by increasing $M$ (more validators) or improving $p_{\max}$ (better validation software), the fork probability can be driven arbitrarily close to zero. This stands in contrast to the prevailing blockchain narrative that treats forks as an acceptable, self-resolving phenomenon through the ``longest chain rule''—a heuristic whose formal guarantees are strictly weaker than the exponential bound we establish.
\end{remark}

%=====================================================================
\section{Theorem 2: 51\% Attack as Effective-M Collapse 51\%攻击作为有效M塌缩定理}
\label{sec:theorem2}
%=====================================================================

The ``51\% attack'' is the canonical threat to blockchain security: an adversary controlling a majority of validation power can rewrite history, double-spend, and censor transactions. We prove that this attack is formally equivalent to \textbf{effective auditor multiplicity collapse}: when $\alpha \geq 0.5$, the effective multiplicity $\Meff$ of honest auditors drops below the consensus threshold $\Mthresh$, making reliable certification impossible.

\subsection{Effective Multiplicity Under Adversarial Capture 对抗捕获下的有效多重性}

\begin{definition}[Effective Honest Multiplicity 有效诚实多重性]
\label{def:Meff_honest}
When fraction $\alpha$ of validators are adversarial, the effective honest multiplicity is:
\begin{equation}
    \Meff(\alpha) = \frac{(1-\alpha)N}{1 + ((1-\alpha)N - 1)\bar{\rho}_h},
    \label{eq:Meff_alpha}
\end{equation}
where $\bar{\rho}_h$ is the average pairwise correlation among honest validators. The adversarial validators are modeled as a single correlated entity (since they coordinate), contributing at most 1 effective auditor regardless of their number.
\end{definition}

\begin{lemma}[Adversarial Block Certification 对抗区块认证]
\label{lem:adv_cert}
Under adversary control fraction $\alpha$, the probability that an invalid block $B_h'$ proposed by the adversary receives $\geq \tau N$ total certifications (honest + adversarial) satisfies:
\begin{equation}
    \Pbb(\text{adv\_block\_certified}) \geq 1 - \exp\left(-2\Meff(\alpha) \cdot \max(0, \alpha - p_{\max})^2\right),
    \label{eq:adv_block_certified}
\end{equation}
where $\Meff(\alpha) = (1-\alpha)N / (1 + ((1-\alpha)N - 1)\bar{\rho}_h)$.
\end{lemma}

\begin{proof}\rigorFull
\textbf{Step 1: Certification structure.} The adversary controls $\alpha N$ validators, all of whom certify the adversarial block. For the block to reach consensus threshold $\tau N$, it needs at least $\tau N - \alpha N = (\tau - \alpha)N$ additional certifications from honest validators.

\textbf{Step 2: Honest certification of invalid block.} Each honest validator $V_j$ independently certifies an invalid block with probability $p_j \leq p_{\max}$. Let $X_h = \sum_{j \in \honestSet} \ind{\hat{v}_j = 1}$ be the number of honest validators who erroneously certify.

The adversarial block is certified iff $\alpha N + X_h \geq \tau N$, i.e., $X_h \geq (\tau - \alpha)N$.

\textbf{Step 3: When $\alpha \geq \tau$.} If the adversary already controls $\geq \tau N$ validators, no honest certifications are needed. The block is certified with probability 1. The effective honest multiplicity is irrelevant because the adversary has captured the consensus mechanism outright.

\textbf{Step 4: When $\alpha < \tau$.} The adversary needs $X_h \geq (\tau - \alpha)N$ honest certifications. The expected number is $\E[X_h] = \sum_{j \in \honestSet} p_j \leq (1-\alpha)N p_{\max}$. The required excess is $(\tau - \alpha)N - (1-\alpha)N p_{\max} = N(\tau - \alpha - (1-\alpha)p_{\max})$.

For the adversary to succeed, $X_h$ must exceed its expectation by at least this amount. By Hoeffding's inequality with effective multiplicity $\Meff(\alpha)$:
\begin{equation}
    \Pbb(X_h \geq (\tau - \alpha)N) \leq \exp\left(-2\Meff(\alpha) \cdot (\tau - \alpha - (1-\alpha)p_{\max})^2 / (1-\alpha)^2\right).
\end{equation}

\textbf{Step 5: Adversarial success lower bound.} The adversary can always choose to propose blocks when network conditions favor certification. Taking the complementary perspective, the probability that the adversary's block is NOT certified is bounded above by the Hoeffding tail, yielding Eq.~\eqref{eq:adv_block_certified}. For large $\alpha$, the exponent becomes small or negative, and the bound becomes vacuous—exactly when the adversary has de facto control. $\square$
\end{proof}

\begin{theorem}[51\% Attack as Effective-M Collapse 51\%攻击作为有效M塌缩定理]
\label{thm:fiftyone}
\rigorFull
Under Assumptions~\ref{ass:A1}--\ref{ass:A3} and~\ref{ass:A7}, there exists a critical adversary fraction $\alpha^*$ such that for all $\alpha > \alpha^*$, the effective honest multiplicity $\Meff(\alpha)$ falls below the consensus threshold $\Mthresh$:
\begin{equation}
    \Meff(\alpha) < \Mthresh \equiv \frac{\ln(2)}{2(\tau - p_{\max})^2},
    \label{eq:Meff_threshold}
\end{equation}
and the probability of detecting an adversarial chain rewrite falls below any fixed confidence level. The critical fraction is:
\begin{equation}
    \alpha^* = \min\left\{\tau,\; \frac{1}{2}\right\}.
    \label{eq:alpha_star}
\end{equation}
Specifically:
\begin{enumerate}[label=(\roman*)]
    \item For $\alpha \geq \tau$: the adversary can certify arbitrary blocks unilaterally. Consensus guarantee is \textbf{void}.
    \item For $\alpha \in [0.5, \tau)$: the adversary can cause forks with probability approaching 1 as $H \to \infty$, and can rewrite history of depth $d$ with expected success after $\sim 2^d$ attempts.
    \item For $\alpha < 0.5$: honest validators form a majority, Theorem~1 applies, and fork probability is exponentially bounded.
\end{enumerate}
\end{theorem}

\begin{proof}\rigorFull
\textbf{Case 1: $\alpha \geq \tau$.} The adversary controls at least $\tau N$ validators. For any block the adversary proposes, all $\alpha N$ adversarial validators certify it, immediately satisfying the consensus condition $\geq \tau N$. The adversary can certify any sequence of blocks, effectively rewriting the entire chain. The honest validators' certifications are irrelevant—their effective multiplicity $\Meff(\alpha)$ does not enter the certification condition because the adversary has already met the threshold unilaterally.

\textbf{Case 2: $\alpha \in [0.5, \tau)$.} The adversary needs $(\tau - \alpha)N$ honest certifications per block. The expected number of honest certifications for an invalid block is $(1-\alpha)N p_{\max}$. The gap is:
\begin{equation}
    \Delta_{\text{adv}} = (\tau - \alpha) - (1-\alpha)p_{\max}.
\end{equation}
For $\alpha \in [0.5, \tau)$, we have $\tau - \alpha \leq \tau - 0.5$. With typical BFT thresholds $\tau \approx 2/3$, this gives $\Delta_{\text{adv}} \leq 1/6 - (1-\alpha)p_{\max}$, which is small.

The probability that honest validators provide sufficient certifications for any single adversarial block is, by Lemma~\ref{lem:adv_cert}:
\begin{equation}
    \Pbb(\text{single\_block\_success}) \leq \exp\left(-2\Meff(\alpha) \cdot \Delta_{\text{adv}}^2\right).
\end{equation}

Over $H$ block proposals, the adversary's probability of succeeding at least once is:
\begin{equation}
    \Pbb(\text{chain\_rewrite}) \geq 1 - \left(1 - \exp(-2\Meff(\alpha) \cdot \Delta_{\text{adv}}^2)\right)^H.
\end{equation}

For $H \gg \exp(2\Meff(\alpha) \cdot \Delta_{\text{adv}}^2)$, this probability approaches 1. Since $\Meff(\alpha)$ decreases as $\alpha$ increases (fewer honest validators), the adversary's success probability grows with $\alpha$.

\textbf{Case 3: $\alpha < 0.5$.} Honest validators form a strict majority. The effective multiplicity is $\Meff(\alpha) = (1-\alpha)N / (1 + ((1-\alpha)N-1)\bar{\rho}_h) > N/2(1+N\bar{\rho}_h)$. Theorem~1 applies directly, giving exponentially small fork probability.

\textbf{Derivation of $\Mthresh$.} The threshold effective multiplicity below which the fork probability exceeds $\varepsilon = 0.5$ is obtained by solving $\exp(-2\Mthresh \cdot (\tau - p_{\max})^2) = 0.5$, yielding $\Mthresh = \ln(2) / (2(\tau - p_{\max})^2)$. This is the minimal effective multiplicity for meaningful consensus—below this, random chance produces forks with probability $> 0.5$ per block.

\textbf{Equivalence to 51\%.} For $\alpha = 0.5$, $\Meff(0.5) = 0.5N / (1 + (0.5N-1)\bar{\rho}_h)$. For large $N$ and small $\bar{\rho}_h$, $\Meff(0.5) \approx 0.5N / 0.5N\bar{\rho}_h = 1/\bar{\rho}_h$. When $\bar{\rho}_h > 0$, this is bounded, and for typical values $\bar{\rho}_h \approx 0.05$, we get $\Meff(0.5) \approx 20$. Comparing with $\Mthresh$: for $\tau = 2/3$, $p_{\max} = 0.01$, $\Mthresh = \ln(2) / (2 \cdot (0.6567)^2) \approx 0.80$, so $\Meff(0.5) > \Mthresh$ and consensus still holds. However, this is only for a single block—the adversary's ability to attempt the attack repeatedly (grinding) over many blocks eventually succeeds, as shown in Case 2. The $50\%$ boundary is thus a \textit{security parameter}, not a sharp phase transition: at exactly $\alpha = 0.5$, security degrades continuously rather than collapsing instantly. $\square$
\end{proof}

\begin{remark}[Beyond 51\%: The Role of $\bar{\rho}$]
The effective multiplicity formula reveals that \textbf{validator correlation} $\bar{\rho}$ is as important as validator count $N$. If all validators run identical software on identical infrastructure ($\bar{\rho} \to 1$), then $\Meff \to 1$ regardless of $N$—a single software bug can compromise the entire validator set. This is the \SCX{} insight that \textbf{auditor diversity is a security resource}: heterogeneous validator implementations (different clients, languages, operating systems) reduce $\bar{\rho}$ and increase $\Meff$, strengthening consensus guarantees without increasing $N$.
\end{remark}

%=====================================================================
\section{Theorem 3: PoS Slashing as Audit Penalty 权益证明罚没作为审计惩罚定理}
\label{sec:theorem3}
%=====================================================================

Proof-of-stake systems deter equivocation through \textbf{slashing}: validators who certify conflicting blocks forfeit their staked tokens. We prove that slashing is structurally identical to the \SCX{} audit penalty $\kappa$, and we derive the \Yajie{} equilibrium condition under which honest validation strictly dominates equivocation.

\subsection{The Slashing Game 罚没博弈}

\begin{definition}[PoS Slashing as Audit Penalty 权益证明罚没作为审计惩罚]
\label{def:slashing_scx}
In an \SCX{}-audited blockchain, each validator $V_j$ posts stake $s_j > 0$. The \textbf{slashing penalty} $\slashPenalty_j = \gamma \cdot s_j$ (with $\gamma \in (0, 1]$) is the amount forfeited upon detected equivocation. This is precisely the \SCX{} audit penalty $\kappa$: the cost incurred when the multi-expert auditor panel detects a false certification.
\end{definition}

\begin{definition}[Validator Strategy Space 验证者策略空间]
\label{def:strategy}
Each validator $V_j$ chooses a strategy:
\begin{itemize}
    \item \textbf{Honest} ($H$): Certify exactly one block per height—the first valid block received. Never equivocate.
    \item \textbf{Equivocate} ($E$): Certify multiple conflicting blocks at the same height, collecting rewards from each while risking slashing.
    \item \textbf{Abstain} ($A$): Certify no blocks, collecting no rewards and risking no slashing.
\end{itemize}
\end{definition}

\begin{lemma}[Equivocation Detection Probability  equivocation检测概率]
\label{lem:equivocation_detection}
Under Assumption~\ref{ass:A5}, if validator $V_j$ equivocates by signing two conflicting blocks $B_h \neq B_h'$, the equivocation is detected with probability:
\begin{equation}
    \Pbb(\text{detection} \mid \text{equivocation}) = 1 - \prod_{k \neq j} (1 - q_k),
    \label{eq:equiv_detection}
\end{equation}
where $q_k$ is the probability that validator $V_k$ observes both conflicting signatures. In a synchronous network (Assumption~\ref{ass:A7}) with gossip propagation, $q_k \to 1$ for all honest validators, so $\Pbb(\text{detection}) \to 1$.
\end{lemma}

\begin{proof}\rigorFull
\textbf{Step 1: Cryptographic detectability.} Equivocation produces two signatures $\sigma_j(B_h)$ and $\sigma_j(B_h')$ by the same validator on conflicting blocks. Any party possessing both signatures can verify that the same validator signed two blocks at the same height, which is a protocol violation. The evidence is cryptographically irrefutable.

\textbf{Step 2: Propagation.} Under Assumption~\ref{ass:A7} (synchronous network with bounded delay $\Delta$), both conflicting blocks propagate to all honest validators within time $\Delta$. Each honest validator $V_k$ who receives both blocks detects the equivocation.

\textbf{Step 3: Detection probability.} An equivocation goes undetected only if no honest validator receives both conflicting blocks. Since adversarial validators may suppress evidence, we consider only honest validators. The probability that honest validator $V_k$ does NOT receive both blocks is at most the probability of a network partition lasting longer than the block propagation time, which is negligible under Assumption~\ref{ass:A7}. Thus $q_k \approx 1$ for all honest $k$, and $\Pbb(\text{detection}) \approx 1$.

\textbf{Step 4: Worst-case bound.} Even under adversarial network conditions, the adversary must prevent ALL honest validators from seeing both blocks. With $|\honestSet|$ honest validators independently connected:
\begin{equation}
    \Pbb(\text{detection}) \geq 1 - \varepsilon^{|\honestSet|},
\end{equation}
where $\varepsilon$ is the per-validator probability of isolation. For any $\varepsilon < 1$, detection is asymptotically certain as $|\honestSet| \to \infty$. $\square$
\end{proof}

\begin{theorem}[PoS Slashing Equilibrium 权益证明罚没均衡定理]
\label{thm:slashing}
\rigorFull
Under Assumptions~\ref{ass:A1}--\ref{ass:A5} and~\ref{ass:A7}, there exists a critical slashing ratio $\gamma^*$ such that for all $\gamma > \gamma^*$, the honest strategy strictly dominates equivocation in the \Yajie{} Nash-Pareto Equilibrium:
\begin{equation}
    U_j(H; \bm{s}_{-j}) > U_j(E; \bm{s}_{-j}), \quad \forall \bm{s}_{-j}.
    \label{eq:honest_dominance}
\end{equation}
The critical slashing ratio is:
\begin{equation}
    \gamma^* = \frac{R_{\text{extra}}}{\Pbb(\text{detection}) \cdot s_j},
    \label{eq:gamma_star}
\end{equation}
where $R_{\text{extra}}$ is the additional reward from equivocating (collecting rewards from multiple blocks). When $\Pbb(\text{detection}) \to 1$, we have $\gamma^* = R_{\text{extra}} / s_j$, meaning the slashed amount must exceed the extra reward.
\end{theorem}

\begin{proof}\rigorFull
\textbf{Step 1: Payoff decomposition.} For honest strategy $H$: validator $V_j$ certifies exactly one block per height, receiving reward $R$ with probability $p_{\text{consensus}}$ (the probability that the certified block achieves consensus) and never facing slashing:
\begin{equation}
    U_j(H) = R \cdot p_{\text{consensus}}.
    \label{eq:U_honest}
\end{equation}

For equivocation strategy $E$: validator $V_j$ certifies $K \geq 2$ conflicting blocks, receiving expected reward $R \cdot \sum_{k=1}^{K} p_{\text{consensus}}^{(k)}$ (sum of consensus probabilities for each certified block), but facing slashing $\slashPenalty_j = \gamma s_j$ with probability $\Pbb(\text{detection})$:
\begin{equation}
    U_j(E) = R \cdot \sum_{k=1}^{K} p_{\text{consensus}}^{(k)} - \gamma s_j \cdot \Pbb(\text{detection}).
    \label{eq:U_equivocate}
\end{equation}

\textbf{Step 2: Reward difference.} The extra reward from equivocation is:
\begin{equation}
    R_{\text{extra}} = R \cdot \left(\sum_{k=1}^{K} p_{\text{consensus}}^{(k)} - p_{\text{consensus}}\right) \leq R \cdot (K - 1),
\end{equation}
since $p_{\text{consensus}}^{(k)} \leq 1$ for all $k$ and $p_{\text{consensus}} \geq 0$. The worst case for the protocol is when all $K$ blocks achieve consensus ($p_{\text{consensus}}^{(k)} \approx 1$), yielding $R_{\text{extra}} \approx R(K-1)$.

\textbf{Step 3: Dominance condition.} Honest strategy dominates equivocation when $U_j(H) > U_j(E)$:
\begin{align}
    R \cdot p_{\text{consensus}} &> R \cdot \sum_{k=1}^{K} p_{\text{consensus}}^{(k)} - \gamma s_j \cdot \Pbb(\text{detection}) \\
    \gamma s_j \cdot \Pbb(\text{detection}) &> R_{\text{extra}} \\
    \gamma &> \frac{R_{\text{extra}}}{s_j \cdot \Pbb(\text{detection})}.
\end{align}

\textbf{Step 4: Best-case for protocol.} When $\Pbb(\text{detection}) = 1$ (guaranteed equivocation detection), the condition simplifies to $\gamma > R_{\text{extra}} / s_j$. With typical parameters: $R_{\text{extra}} \approx 2R$ (equivocating on two blocks), $s_j$ being the validator's stake, the slashing penalty must exceed twice the per-block reward relative to stake. For Ethereum PoS, validators stake 32 ETH and earn $\sim 0.00003$ ETH per attestation, giving $\gamma^* \approx 2 \times 10^{-6}$—a trivially satisfied condition. The binding constraint is not $\gamma$ but the validator's total stake at risk.

\textbf{Step 5: \Yajie{} NPE.} The \Yajie{} Nash-Pareto Equilibrium requires that the equilibrium strategy profile is both a Nash equilibrium and Pareto efficient among the honest validators. The honest strategy profile $\bm{s}^H = (H, H, \ldots, H)$ is:
\begin{itemize}
    \item \textbf{Nash}: By the dominance condition, no validator can unilaterally improve by deviating to $E$ when $\gamma > \gamma^*$.
    \item \textbf{Pareto efficient}: All honest validators receive $U_j(H) = R \cdot p_{\text{consensus}}$. Any deviation to $E$ introduces slashing risk for the deviator and potentially reduces $p_{\text{consensus}}$ for others (by creating forks), making others worse off. Thus $\bm{s}^H$ is Pareto efficient among honest validators.
\end{itemize}
This establishes $\bm{s}^H$ as a \Yajie{} NPE. $\square$
\end{proof}

\begin{corollary}[Slashing as the Invisible Hand 罚没作为看不见的手]
\label{cor:invisible_hand}
The slashing mechanism implements \SCX{} audit incentives without centralized enforcement. The protocol itself is the auditor: any validator can submit evidence of equivocation and receive a portion of the slashed stake (whistleblower reward), creating a self-policing equilibrium. This is structurally identical to the \SCX{} framework where the auditor community collectively enforces the audit penalty $\kappa$, with the crucial difference that blockchain automates enforcement via smart contracts rather than relying on institutional action.
\end{corollary}

\begin{corollary}[Slashing-Audit Isomorphism 罚没-审计同构]
\label{cor:isomorphism}
The mapping between PoS slashing and \SCX{} audit is an isomorphism:
\begin{center}
\begin{tabular}{l|l}
\toprule
\textbf{PoS Slashing} & \textbf{\SCX{} Audit} \\
\midrule
Stake $s_j$ & Audit bond / reputation at risk \\
Slashing penalty $\gamma s_j$ & Audit penalty $\kappa$ \\
Equivocation evidence & Detected data fabrication \\
Whistleblower reward & Community verification incentive \\
Protocol-enforced penalty & \Spring{} permanent record + community action \\
\bottomrule
\end{tabular}
\end{center}
This isomorphism means that \textbf{all \SCX{} theorems about audit optimality apply directly to PoS consensus}, and conversely, blockchain slashing mechanisms can inform the design of \SCX{} audit penalties in other domains (scientific publishing, governance, journalism).
\end{corollary}

%=====================================================================
\section{Hash Chain as \Spring's Permanent Memory 哈希链作为春耕永久记忆}
\label{sec:spring}
%=====================================================================

\subsection{The Immutable Audit Trail 不可变审计轨迹}

The \Spring{} mechanism in \SCX{} provides \textbf{permanent memory}: once an audit event is recorded, it cannot be erased or modified. The blockchain's hash chain implements this property natively and with cryptographic strength.

\begin{theorem}[Hash Chain = \Spring{} $\Mt$ 哈希链即春耕永久记忆]
\label{thm:spring_hash}
\rigorFull
Under Assumption~\ref{ass:A6} (hash preimage resistance), the blockchain hash chain $\chain = (B_0, \ldots, B_H)$ with $B_h.\text{parent\_hash} = \SHA256(B_{h-1})$ implements \Spring's permanent memory $\Mt$: the sequence of validator multiplicities $\{M_h\}_{h=1}^{H}$ and their certifications are permanently and immutably recorded. Specifically:
\begin{enumerate}[label=(\roman*)]
    \item \textbf{Immutability}: Modifying any block $B_h$ changes its hash, which changes $B_{h+1}.\text{parent\_hash}$, requiring modification of all subsequent blocks—computationally infeasible under preimage resistance.
    \item \textbf{Verifiability}: Any party can verify the entire audit trail by recomputing hashes, confirming that the recorded $\{M_h\}$ has not been altered.
    \item \textbf{Permanence}: The record persists as long as any copy of the chain exists. Unlike institutional memory (which can be lost, destroyed, or manipulated), cryptographic memory is self-verifying and replicable.
\end{enumerate}
\end{theorem}

\begin{proof}\rigorFull
\textbf{(i) Immutability.} Suppose an adversary modifies block $B_h$ to $B_h'$, changing the recorded validator set or certifications. The new hash $h' = \SHA256(B_h')$ differs from the original $h = \SHA256(B_h)$ (by collision resistance of SHA-256). Block $B_{h+1}$ contains $h$ in its parent\_hash field. To maintain chain consistency, the adversary must either: (a) modify $B_{h+1}$ to reference $h'$, changing $B_{h+1}$'s hash and cascading the modification through all subsequent blocks; or (b) find $B_h' \neq B_h$ with $\SHA256(B_h') = \SHA256(B_h)$, violating preimage resistance (Assumption~\ref{ass:A6}). Option (a) requires $O(2^{256})$ work for the final block's proof-of-work (in PoW) or control of $>2/3$ of stake (in PoS) at every subsequent height—both infeasible under standard security assumptions.

\textbf{(ii) Verifiability.} Given the chain, any verifier recomputes $h_1 = \SHA256(B_1)$, checks that $h_1$ matches $B_2.\text{parent\_hash}$, and iterates for all blocks. If all hashes match, the chain is internally consistent. The verifier can then extract $\{M_h\}$ from block headers and confirm the recorded audit multiplicity at each height.

\textbf{(iii) Permanence.} The chain is replicated across all full nodes. Destroying the record requires simultaneously destroying all copies, which is infeasible for any sufficiently decentralized network. The \Spring{} memory is thus \textbf{information-theoretically permanent} modulo the continued existence of at least one honest full node. $\square$
\end{proof}

\subsection{$\Mt$ as Dynamic Security Parameter 动态安全参数}

The \Spring{} record $\Mt = (M_1, M_2, \ldots, M_H)$ is not merely a historical curiosity—it is a \textbf{dynamic security parameter} that enables real-time auditing of consensus health.

\begin{definition}[$\Mt$ Trajectory Analysis $\Mt$轨迹分析]
\label{def:Mt_trajectory}
The sequence $\Mt$ reveals:
\begin{enumerate}[label=(\roman*)]
    \item \textbf{Validator participation trends}: Declining $\Mt$ signals validator exodus, potentially preceding a security crisis. A sudden drop in $M_h$ at height $h$ may indicate a coordinated validator withdrawal (e.g., in anticipation of a hard fork or regulatory action).
    \item \textbf{Fork events}: Spikes in $M_h$ variance across competing chain tips indicate consensus instability—exactly the fork condition analyzed in Theorem~\ref{thm:fork}.
    \item \textbf{Censorship detection}: If certain transaction types consistently appear in blocks with lower $M_h$, this suggests validators are selectively certifying based on transaction content.
    \item \textbf{Regime shifts}: \Spring{} gating detects transitions from ``normal consensus'' ($M_h$ stable, high) to ``degraded consensus'' ($M_h$ declining, volatile) using the same change-point detection methodology applied in \SCX{} governance auditing.
\end{enumerate}
\end{definition}

\begin{corollary}[\Spring{} Gating for Blockchain Security 春耕门控区块链安全]
\label{cor:spring_gating}
Define the security regime indicator:
\begin{equation}
    R_t = \ind{\Meff(t) \geq \Mthresh \text{ and } \Var(M_{t-W:t}) \leq \sigma^2_{\max}},
    \label{eq:regime}
\end{equation}
where $W$ is a rolling window and $\sigma^2_{\max}$ is the maximum tolerable variance. A transition $R_t: 1 \to 0$ triggers a \Spring{} gate: the system enters a heightened security mode (longer confirmation times, increased slashing penalties, or validator set rotation). This is the blockchain analogue of \Spring{} gating in governance: when the statistical properties of the audit trail signal regime degradation, automatic defensive measures engage.
\end{corollary}

%=====================================================================
\section{Nakamoto Consensus as Implicit M-Declaration 中本聪共识作为隐式M声明}
\label{sec:nakamoto}
%=====================================================================

\subsection{The Implicit Guarantee 隐式保证}

Nakamoto consensus~\cite{nakamoto2008}—the mechanism underlying Bitcoin and early proof-of-work systems—achieves distributed agreement through proof-of-work and the longest-chain rule. We show that this constitutes an \textbf{implicit M-declaration} without formal guarantee.

\begin{definition}[Nakamoto Consensus 中本聪共识]
\label{def:nakamoto}
In Nakamoto consensus:
\begin{enumerate}[label=(\roman*)]
    \item Miners compete to solve a proof-of-work puzzle; the winner proposes the next block.
    \item Other miners implicitly ``certify'' the block by building on top of it (extending the chain).
    \item The longest valid chain is considered the authoritative ledger state.
    \item There is no explicit validator set, no explicit certification step, and no explicit $M$-declaration.
\end{enumerate}
\end{definition}

\begin{theorem}[Nakamoto as Implicit M-Declaration 中本聪作为隐式M声明]
\label{thm:nakamoto_implicit}
In Nakamoto consensus, each block $B_h$ carries an \textbf{implicit} M-declaration: the number of miners who have built on top of $B_h$ serves as a proxy for $M_h$. However, this implicit M lacks formal guarantees:
\begin{enumerate}[label=(\roman*)]
    \item \textbf{No explicit certification}: Miners building on $B_h$ do not explicitly verify its contents—they may extend a block without validating all transactions, relying on the economic incentive that invalid blocks will be orphaned. This is the \textbf{SPV mining problem}.
    \item \textbf{Delayed confirmation}: The implicit $M_h$ is only knowable retrospectively, after $k$ confirmations. At time of block proposal, $M_h = 1$ (only the proposing miner has seen it). The effective $M$ grows with confirmations: $M_h(k) = k$.
    \item \textbf{No independence guarantee}: All miners extending the chain share the same software implementation (Bitcoin Core), creating $\bar{\rho} \approx 1$ for implementation-specific bugs. The effective multiplicity is $\Meff \approx 1$ regardless of $k$.
    \item \textbf{Probabilistic, not absolute}: The Nakamoto guarantee is that after $k$ confirmations, the probability of reversal decays as $O(\exp(-ck))$ for some constant $c$ depending on adversarial hash power. This is a \textit{probabilistic} guarantee, not an \textit{absolute} certification—there is always a non-zero probability of chain reorganization.
\end{enumerate}
\end{theorem}

\begin{proof}\rigorFull
\textbf{(i) SPV mining.} Simplified Payment Verification (SPV) mining allows miners to extend a block without validating its full transaction contents, verifying only the block header. An SPV miner implicitly certifies block structure but not transaction validity. The effective certification is degraded: $\hat{v}_j(B_h)$ may equal 1 even when $B_h$ contains invalid transactions, if those transactions are not checked. This inflates the apparent $M$ without providing genuine audit value.

\textbf{(ii) Retrospective M.} Let $C_h(t)$ be the number of blocks built on top of $B_h$ at time $t$. At $t = 0$ (block proposal), $C_h(0) = 0$, so $M_h(0) = 1$ (the proposer alone). After $k$ blocks, $M_h = k+1$. The \SCX{} framework requires $M$ to be declared \textit{at the time of certification}; the retrospective nature of Nakamoto $M$ means the system operates with $M=1$ during the critical initial period when forks are most likely.

\textbf{(iii) Software monoculture.} Over 95\% of Bitcoin nodes run Bitcoin Core. A consensus bug in Bitcoin Core affects essentially all miners simultaneously, giving $\bar{\rho} \approx 1$ and $\Meff \approx 1$ for software-induced faults. This is the exact vulnerability that Theorem~2 identifies: when $\Meff$ collapses, consensus guarantees vanish regardless of nominal $N$.

\textbf{(iv) Probabilistic guarantee.} The Nakamoto security analysis~\cite{rosenfeld2014} shows that a transaction with $k$ confirmations is reversed with probability:
\begin{equation}
    \Pbb(\text{reversal after } k) \leq \left(\frac{\alpha}{1-\alpha}\right)^k,
    \label{eq:nakamoto_reversal}
\end{equation}
where $\alpha$ is the adversary's hash power fraction. For $\alpha \to 0.5$, this bound approaches 1 for any finite $k$. In contrast, the \SCX{} bound (Theorem~\ref{thm:fork}) is $\exp(-2M(\tau - p_{\max})^2)$, which is exponentially tight regardless of $\alpha$ (as long as $\alpha < 0.5$ and $M$ is explicitly declared).

\textbf{The structural gap.} Nakamoto consensus replaces explicit M-declaration with a \textbf{market mechanism}: the cost of proof-of-work is assumed to align miner incentives with chain validity. This alignment is an economic argument, not a mathematical proof. The \SCX{} framework requires mathematical proof: $M$ must be declared, validators must be independent, and certification must be explicit. The gap between ``economic incentive'' and ``mathematical guarantee'' is the gap between Nakamoto consensus and \SCX{}-audited consensus. $\square$
\end{proof}

\begin{remark}[The Irony of Bitcoin's Security]
Bitcoin is often described as ``the most secure blockchain'' because of its immense hash power. Under \SCX{} analysis, Bitcoin's security derives not from hash power per se, but from the \textbf{implicit auditor multiplicity} created by the network of full nodes that independently verify every block. A Bitcoin full node is an \SCX{} auditor: it validates all transactions, checks consensus rules, and rejects invalid blocks. The true $M$ for Bitcoin is not the number of miners but the number of \textbf{full nodes}—and this $M$ is never explicitly declared, varies over time, and includes many nodes running identical software.
\end{remark}

%=====================================================================
\section{Consensus Mechanism Analysis 共识机制分析}
\label{sec:mechanisms}
%=====================================================================

We apply the \SCX{} audit framework to four major consensus mechanisms, computing their effective $M$, fork bounds, and slashing properties.

\subsection{Bitcoin (PoW Nakamoto Consensus)}

\begin{itemize}
    \item \textbf{Validator model}: Miners propose blocks; implicit certification by subsequent miners building on the chain.
    \item \textbf{Explicit $M$}: None. Implicit $M =$ number of confirmations $+ 1$.
    \item \textbf{$\Meff$}: $\Meff \approx 1$ for software bugs ($\bar{\rho} \approx 1$); $\Meff \approx k$ for hash-power-based security after $k$ confirmations.
    \item \textbf{Fork bound}: $\Pbb(\text{fork}) \leq (\alpha/(1-\alpha))^k$, where $\alpha$ is adversarial hash power (weaker than Theorem~\ref{thm:fork}).
    \item \textbf{Slashing}: None. The penalty for malicious mining is wasted electricity (implicit $\kappa =$ electricity cost).
    \item \textbf{\SCX{} assessment}: Bitcoin achieves security through economic incentives, not formal certification. The lack of explicit M-declaration means security is always retrospective and probabilistic. However, the full node network provides a de facto auditor community whose size is the actual security parameter—but this parameter is unmeasured and undeclared.
\end{itemize}

\subsection{Ethereum Proof-of-Stake (Gasper)}

\begin{itemize}
    \item \textbf{Validator model}: Explicit validator set ($\sim 1$M validators). Each validator attests to one block per slot.
    \item \textbf{Explicit $M$}: Number of attestations per block (typically $\sim 500$K per epoch).
    \item \textbf{$\Meff$}: High due to large $N$, but reduced by client monoculture ($\sim 70\%$ Prysm/Lighthouse). Effective $\Meff \approx N / (1 + N \cdot 0.3) \approx 3$ for implementation bugs.
    \item \textbf{Fork bound}: Theorem~\ref{thm:fork} applies with $M \approx 500$K, $\tau = 2/3$, $p_{\max} \approx 0.001$ (attestation errors). $\Pbb(\text{fork}) \leq 2\exp(-2 \cdot 500\text{K} \cdot 0.44^2) \approx 10^{-84000}$, effectively zero.
    \item \textbf{Slashing}: Explicit. $\gamma = 1/32$ of stake for equivocation (plus additional penalties for surround voting). Theorem~\ref{thm:slashing} confirms honest dominance for any $\gamma > 0$.
    \item \textbf{\SCX{} assessment}: Ethereum PoS is the closest existing system to \SCX{}-audited consensus. It has explicit validators, explicit attestations ($M$), and explicit slashing ($\kappa$). The main weakness is client monoculture ($\bar{\rho}$ too high) and the complexity of the fork-choice rule (Gasper), which introduces additional failure modes beyond the simple fork model.
\end{itemize}

\subsection{Tendermint (BFT)}

\begin{itemize}
    \item \textbf{Validator model}: Fixed validator set (typically $N = 100$--$150$). Two-thirds supermajority required per block.
    \item \textbf{Explicit $M$}: Number of prevotes/precommits per block ($\sim 100$).
    \item \textbf{$\Meff$}: $M \approx 100$, $\bar{\rho}$ depends on validator diversity (typically low if validators are independent organizations).
    \item \textbf{Fork bound}: Theorem~\ref{thm:fork} with $M = 100$, $\tau = 2/3$, $p_{\max} = 0$. $\Pbb(\text{fork}) \leq 2\exp(-2 \cdot 100 \cdot (2/3)^2) \approx 2\exp(-88.9) \approx 5 \times 10^{-39}$. Tendermint provides \textbf{absolute finality}: once a block receives $>2/3$ prevotes, no fork is possible (assuming $<1/3$ Byzantine validators).
    \item \textbf{Slashing}: Explicit. $\gamma = 5\%$ of stake for equivocation (light client attack: full slashing).
    \item \textbf{\SCX{} assessment}: Tendermint is the gold standard for \SCX{}-audited consensus: explicit $M$, explicit $\kappa$, absolute finality. The limitation is scalability: $N$ is bounded by communication complexity ($O(N^2)$ messages per block).
\end{itemize}

\subsection{Solana (Proof-of-History + Tower BFT)}

\begin{itemize}
    \item \textbf{Validator model}: $\sim 2000$ validators. Tower BFT uses PoH as a global clock; validators vote on forks.
    \item \textbf{Explicit $M$}: Number of votes per block (variable).
    \item \textbf{$\Meff$}: Large $N$ but high hardware requirements create centralization pressure ($\bar{\rho}$ elevated due to shared infrastructure dependencies).
    \item \textbf{Fork bound}: Theorem~\ref{thm:fork} applies but with caveat: PoH clock synchronization failures can create forks that are not captured by the simple independent-validator model.
    \item \textbf{Slashing}: Present but rarely enforced. Slashing is not automatic—requires manual governance intervention, reducing $\Pbb(\text{detection})$ and weakening Theorem~\ref{thm:slashing}'s dominance condition.
    \item \textbf{\SCX{} assessment}: Solana's architecture introduces complex dependencies (PoH clock) that violate the independence assumption (Assumption~\ref{ass:A2}). The effective $\Meff$ for clock-failure-induced forks is $\Meff \approx 1$ (all validators share the same clock). This explains Solana's historical susceptibility to consensus failures despite large nominal $N$.
\end{itemize}

%=====================================================================
\section{Discussion 讨论}
\label{sec:discussion}
%=====================================================================

\subsection{Summary of Results 结果总结}

We have established three fundamental theorems connecting blockchain consensus to the \SCX{} multi-expert audit framework:

\begin{enumerate}[label=\textbf{T\arabic*}:, leftmargin=*]
    \item \textbf{Fork Boundary}: The probability of consensus failure (fork) is bounded by $\exp(-2\Meff \cdot (\tau - p_{\max})^2)$, making forks exponentially rare as validator multiplicity grows. Forks are not a feature to be resolved by longest-chain tiebreakers—they are audit failures to be prevented by sufficient $M$.
    \item \textbf{51\% = $\Meff$ Collapse}: The 51\% attack is formally equivalent to effective auditor multiplicity dropping below the certification threshold. The critical parameter is not raw hash power or stake but $\Meff$, which depends on both validator count $N$ and validator correlation $\bar{\rho}$.
    \item \textbf{Slashing = Audit Penalty}: PoS slashing implements the \SCX{} audit penalty $\kappa$ with automated enforcement. The \Yajie{} equilibrium condition $\gamma > R_{\text{extra}} / (s_j \cdot \Pbb(\text{detection}))$ is easily satisfied in practice, explaining why PoS systems with slashing have not experienced economically rational equivocation.
\end{enumerate}

\subsection{The Structural Gap: Implicit vs.\ Explicit M 隐式与显式M的结构性差距}

The central insight of this paper is the distinction between \textbf{implicit} and \textbf{explicit} M-declaration:

\begin{itemize}
    \item \textbf{Implicit M} (Nakamoto consensus, Bitcoin): The number of validators is inferred from blockchain data (confirmations, hash power). No validator explicitly declares ``I have verified this block.'' The security guarantee is probabilistic, retrospective, and contingent on economic assumptions.
    \item \textbf{Explicit M} (Tendermint, Ethereum PoS with attestations): Each validator explicitly signs the block, providing cryptographically verifiable evidence of certification. The security guarantee is absolute (for BFT), prospective (at block time), and mathematically provable.
\end{itemize}

The \SCX{} framework requires explicit M. The trend in blockchain design—from Bitcoin's implicit certification to Ethereum's explicit attestations to Tendermint's explicit prevotes/precommits—is a trajectory toward \SCX{}-compatible consensus. This paper provides the formal vocabulary to describe, analyze, and accelerate this trajectory.

\subsection{$\bar{\rho}$ as the Hidden Vulnerability $\bar{\rho}$作为隐藏漏洞}

The effective multiplicity formula $\Meff = M / (1 + (M-1)\bar{\rho})$ reveals a structural vulnerability that is underappreciated in blockchain security analysis: \textbf{validator correlation} $\bar{\rho}$ can destroy security even when $M$ is large. This has practical implications:

\begin{enumerate}[label=(\roman*)]
    \item \textbf{Client diversity}: Blockchain protocols should mandate or incentivize multiple independent client implementations. Ethereum's client diversity initiative is a direct (if unrecognized) implementation of \SCX{} correlation reduction.
    \item \textbf{Infrastructure diversity}: Validators should use diverse hosting providers, geographic regions, and network topologies to reduce correlated failure modes.
    \item \textbf{Software supply chain}: A single compromised dependency (e.g., a widely-used cryptographic library) creates $\bar{\rho} \to 1$ for all validators using that library—an attack surface that $\Meff$ analysis exposes.
\end{enumerate}

\subsection{Limitations 局限性}

\begin{enumerate}[label=\textbf{L\arabic*}:, leftmargin=*]
    \item \textbf{Synchrony assumption}: Theorem~\ref{thm:fork} assumes synchronous message delivery (Assumption~\ref{ass:A7}). Under partial synchrony or asynchrony, fork probabilities may be higher. Extending the bounds to partially synchronous models is an open problem.
    \item \textbf{Independent errors}: Assumption~\ref{ass:A2} (independent validation) abstracts away systemic risks: a protocol-level bug affects all validators identically, creating $\bar{\rho} = 1$ regardless of implementation diversity. Our framework captures this through $\Meff$ but cannot prevent it.
    \item \textbf{No defense against $>$50\% attacks}: Theorem~\ref{thm:fiftyone} proves that $\alpha > 0.5$ breaks consensus; it does not provide a defense. The ``defense'' is ensuring $\alpha < 0.5$ through economic mechanism design.
    \item \textbf{Scalability vs.\ security}: Larger $M$ improves security (Theorem~\ref{thm:fork}) but reduces throughput (more signatures to verify, more messages to propagate). The optimal $M$ is a trade-off that depends on the application's security requirements.
    \item \textbf{Complex fork-choice rules}: Modern protocols (Ethereum's LMD-GHOST, Solana's Tower BFT) have fork-choice rules whose formal analysis under \SCX{} is more complex than the simple threshold model in Theorem~\ref{thm:fork}. Extending the framework to these rules is future work.
\end{enumerate}

\subsection{The Path to SCX-Audited Consensus 走向SCX审计共识之路}

We envision a future where blockchain consensus protocols are designed from first principles to be \SCX{}-auditable:

\begin{enumerate}[label=(\roman*)]
    \item \textbf{Mandatory M-declaration}: Every block must record the exact set of validators who certified it, with cryptographic signatures. This is already standard practice; our contribution is formalizing \textit{why} it matters.
    \item \textbf{Explicit $\Meff$ computation}: Protocols should compute and publish $\Meff$ accounting for validator correlation, enabling real-time security monitoring.
    \item \textbf{$\bar{\rho}$ minimization}: Protocol design should incentivize validator diversity (multiple clients, diverse infrastructure) to minimize $\bar{\rho}$.
    \item \textbf{Formal $\kappa$ calibration}: Slashing penalties should be calibrated using Theorem~\ref{thm:slashing}, not set arbitrarily.
    \item \textbf{\Spring{} gating}: Automatic security escalation when $\Meff$ drops or $\Mt$ trajectory signals regime degradation.
\end{enumerate}

%=====================================================================
\section{Conclusion 结论}
\label{sec:conclusion}
%=====================================================================

Blockchain consensus and \SCX{} multi-expert audit are two manifestations of the same mathematical structure: distributed certification of claims about state. By formalizing this connection, we have:

\begin{enumerate}[label=(\roman*)]
    \item Provided exponential bounds on consensus failure (fork) probability, showing that forks are audit failures whose likelihood can be driven arbitrarily low by increasing validator multiplicity $M$.
    \item Proven the formal equivalence between 51\% attacks and effective auditor multiplicity collapse, clarifying that security depends on $\Meff$ (which incorporates validator correlation $\bar{\rho}$) rather than raw validator count.
    \item Established the \Yajie{} equilibrium condition for PoS slashing, proving that honest validation strictly dominates equivocation when the slashing penalty exceeds a computable threshold.
    \item Demonstrated that the hash chain implements \Spring's permanent memory $\Mt$, providing an immutable, self-verifying audit trail of consensus health.
    \item Identified Nakamoto consensus as implicit M-declaration without formal guarantee, highlighting the structural gap that explicit \SCX{} certification fills.
\end{enumerate}

The blockchain community has independently converged on many \SCX{} principles—explicit validator sets, slashing penalties, client diversity—without the formal vocabulary to articulate why these design choices matter. This paper provides that vocabulary. The next generation of consensus protocols should be \textbf{SCX-native}: designed from the ground up with explicit M-declaration, $\Meff$ optimization, $\bar{\rho}$ minimization, and formal $\kappa$ calibration as first-class requirements.

\begin{center}
\large\bfseries
No M, No Consensus. 无M，无共识。\\
Declare M explicitly. 显式声明M。\\
Fork = audit failure, bounded by M. 分叉 = 审计失败，由M界定。\\
51\% = $\Meff$ collapse, not hash power. 51\% = $\Meff$塌缩，而非算力。\\
Slashing = $\kappa$, the invisible hand. 罚没 = $\kappa$，看不见的手。\\
Hash chain = \Spring{} memory, permanent. 哈希链 = 春耕记忆，永恒。\\
Nakamoto = implicit M, no formal guarantee. 中本聪 = 隐式M，无形式化保证。\\
Consensus belongs to verifiers. 共识属于验证者。
\end{center}

%=====================================================================
\begin{thebibliography}{99}

\bibitem{SCX2025}
\SCX{} Research Collective.
The \SCX{} Framework: Structured Causal eXamination for Multi-Expert Audit.
Technical report, 2025--2026.

\bibitem{SCX_science_audit}
\SCX{} Research Collective.
The \SCX{} Audit Mandate: Why M-Parameter Declaration Must Be a Prerequisite for Scientific Publication.
Technical report, June 2026.

\bibitem{SCX_governance}
\SCX{} Research Collective.
\SCX{} Audit of Governance: Game-Theoretic Foundations of Transparency Under Multi-Expert Verification.
Technical report, 2026.

\bibitem{nakamoto2008}
S.~Nakamoto.
Bitcoin: A Peer-to-Peer Electronic Cash System.
2008.

\bibitem{buterin2014}
V.~Buterin.
Ethereum: A Next-Generation Smart Contract and Decentralized Application Platform.
2014.

\bibitem{buchman2016}
E.~Buchman.
Tendermint: Byzantine Fault Tolerance in the Age of Blockchains.
Master's thesis, University of Guelph, 2016.

\bibitem{yakovenko2018}
A.~Yakovenko.
Solana: A New Architecture for a High Performance Blockchain.
2018.

\bibitem{hoeffding1963}
W.~Hoeffding.
Probability Inequalities for Sums of Bounded Random Variables.
\emph{Journal of the American Statistical Association}, 58(301):13--30, 1963.

\bibitem{rosenfeld2014}
M.~Rosenfeld.
Analysis of Hashrate-Based Double Spending.
\emph{arXiv preprint arXiv:1402.2009}, 2014.

\bibitem{spence1973}
M.~Spence.
Job Market Signaling.
\emph{Quarterly Journal of Economics}, 87(3):355--374, 1973.

\bibitem{castro1999}
M.~Castro and B.~Liskov.
Practical Byzantine Fault Tolerance.
In \emph{OSDI}, 1999.

\bibitem{garey1979}
M.~R.~Garey and D.~S.~Johnson.
\emph{Computers and Intractability: A Guide to the Theory of NP-Completeness}.
W.~H.~Freeman, 1979.

\bibitem{daian2020}
P.~Daian et al.
Flash Boys 2.0: Frontrunning in Decentralized Exchanges, Miner Extractable Value, and Consensus Instability.
In \emph{IEEE S\&P}, 2020.

\bibitem{neu2022}
J.~Neu et al.
Ebb-and-Flow Protocols: A Resolution of the Availability-Finality Dilemma.
In \emph{IEEE S\&P}, 2022.

\bibitem{gilad2017}
Y.~Gilad et al.
Algorand: Scaling Byzantine Agreements for Cryptocurrencies.
In \emph{SOSP}, 2017.

\end{thebibliography}

\end{document}
